\documentclass[11pt]{article}
\usepackage[margin=1in]{geometry}
\usepackage{hyperref}
\usepackage{booktabs}
\usepackage{natbib}
\usepackage{titlesec}

\title{Verbosity Without Veracity: A Reproducible Multi-Model Probe of\\
Large Language Models on Long-Form Scientific Synthesis}

\author{Fred Zimmerman\\
\small AI Lab for Book-Lovers, Nimble Books LLC\\
\small \texttt{wfz@nimblebooks.com}}

\date{Experiment conducted November 2025; manuscript June 2026}

\begin{document}
\maketitle

\begin{abstract}
We report a small, fully reproducible probe of how a panel of seven large language
models (LLMs) behaves when asked to draft long-form scientific-assessment text
\emph{in the style of} an Intergovernmental Panel on Climate Change (IPCC)
Working Group~II report. Using a single fixed prompt set, identical task definitions,
and a uniform LLM-judge rubric, we generated up to 29 chapters per model
($\sim$330k words in the largest configuration) and scored each for quality and
factual issues. Our central observation is counter-intuitive and policy-relevant:
\textbf{output length was negatively associated with factual reliability.} The most
verbose model produced roughly $2.8\times$ as many flagged fact-check issues as a more
concise, higher-scoring model on matched chapters, while scoring marginally lower
overall. Separately, citation grounding was uniformly weak across all models under
non-citation-enforcing prompts. We frame these results strictly as a
\emph{measurement of model behavior}, not as a proposal to automate scientific
assessment, and we discuss threats to validity---chiefly the single-LLM-judge
design---at length. All prompts, scoring code, and rubrics are released. This is a
negative-leaning, hygiene-oriented contribution intended to caution against
length-as-quality heuristics in AI-for-science pipelines. This work is not
affiliated with or endorsed by the IPCC.
\end{abstract}

\section{Introduction}
LLMs are increasingly proposed as accelerators for scientific synthesis: literature
review, drafting, and summarization. Whether they \emph{should} be used in any
high-stakes assessment setting is a separate normative question we do not address;
the IPCC process is deliberative, expert-governed, and accountable, and nothing here
substitutes for it. Instead we ask a narrow empirical question: \emph{given} the task
of drafting assessment-style text from a published author outline, how do current
models behave, and how do they differ under identical conditions?

This question matters because failure modes in long-form synthesis are not well
characterized relative to short-answer benchmarks. Hallucination and unfaithfulness
are documented for summarization and QA \citep{maynez2020faithfulness, ji2023survey,
lin2021truthfulqa}, and holistic evaluation efforts have mapped model behavior across
many axes \citep{liang2022holistic}. We add a focused, reproducible data point on the
specific regime of multi-thousand-word, citation-bearing scientific prose.

\section{Method}
\paragraph{Task.} Models were prompted with the published IPCC AR7 WGII author outline
and asked to draft chapters (Summary for Policymakers, Technical Summary, and
assessment chapters). ``AR7 / Working Group~II'' names only the \emph{target format}
the models imitate. Some models adopt a first-person ``Coordinating Lead Author''
voice; this is a prompt artifact, not a claim of any role.

\paragraph{Models.} Two tiers, all November-2025 versions. \emph{Premium tier:} seven
models (OpenAI GPT-5, Google Gemini~2.5~Pro, xAI Grok~3, Anthropic Claude~Sonnet~4,
Mistral Mixtral~8x7B, Qwen QwQ-32B, DeepSeek-32B) $\times$ 3 chapters.
\emph{Flash/lite tier:} three models (Gemini~2.5~Flash, Claude~Haiku~4.5,
GPT-4o-mini) $\times$ 29 chapters (full body), which supplies the headline finding.

\paragraph{Evaluation.} A single LLM judge (Gemini~2.5~Pro) applied (i) a multi-dimensional
1--7 Likert rubric (accuracy, style fidelity, uncertainty-language calibration,
synthesis, comprehensiveness, citation quality) and (ii) a fact-check pass categorizing
issues as critical / major / minor. We deliberately note that LLM-as-judge is a biased
instrument \citep{zheng2023judging, panickssery2024llm} and treat its scores as a
relative, reproducible signal rather than ground truth (Section~\ref{sec:threats}).

\section{Results (November 2025 snapshot)}
\subsection{Length did not buy reliability}
On the full 29-chapter run, scored on matched chapters:

\begin{table}[h]
\centering
\begin{tabular}{lrrr}
\toprule
Model & Words & Quality (1--7) & Fact-check issues \\
\midrule
Gemini 2.5 Flash & $\sim$141k & 6.02 & 23 \\
Claude Haiku 4.5 & $\sim$160k & 5.81 & 64 \\
GPT-4o-mini      & $\sim$29k  & 3.19 & 16 (insufficient content) \\
\bottomrule
\end{tabular}
\caption{The most verbose model (Haiku, $+13\%$ words) drew $\sim$2.8$\times$ more
flagged issues while scoring slightly lower. The shortest model failed for the
opposite reason---too little substance to assess.}
\end{table}

A plausible mechanism is a \emph{specificity paradox}: longer prose introduces more
concrete claims (dates, magnitudes) and thus more surface area for unsupported
assertions---consistent with faithfulness findings in summarization
\citep{maynez2020faithfulness}. The practical implication is that word-count or
length-reward is a poor proxy for synthesis quality.

\subsection{Citation grounding was uniformly weak}
Under prompts that did not mandate citations, citation-quality scores were low
($\approx$3.86--4.57/7) for \emph{every} model: fluent, well-structured, style-faithful
prose that was largely untethered to verifiable references. This is the most
policy-relevant failure mode and motivates a citation-enforcing prompt variant
(released for follow-up).

\subsection{Premium tier}
In the 3-chapter premium run the evaluator ranked Gemini~2.5~Pro highest (7.00/7),
with GPT-5 and Grok~3 close behind; smaller open models produced outlines or thin
prose, and two models timed out on the longest chapter---itself a long-context
reliability signal.

\section{Threats to validity}\label{sec:threats}
\textbf{Single LLM judge.} Scores derive from one model acting as judge; LLM judges
exhibit self-preference and verbosity biases \citep{panickssery2024llm,
zheng2023judging}. Scores are relative, not authoritative.
\textbf{No human expert validation.} No climate scientist reviewed any output.
\textbf{Snapshot.} Results are Nov-2025 model versions and will not replicate on
current models. \textbf{Small premium $n$.} Three chapters per model is directional.
\textbf{Prompt sensitivity.} Results may shift with different prompts; ours are public.

\section{Conclusion}
We contribute a reproducible harness and one robust, slightly counter-intuitive
observation: under our conditions, verbosity correlated with \emph{more} factual
errors, and citation grounding was uniformly weak. We make no claim that AI can or
should write scientific assessments; we claim only that, if asked, current models
fail in measurable, characterizable ways. We release all prompts, code, and rubrics
and invite alternative evaluators and human spot-checks.

\paragraph{Data and code.}
\url{https://github.com/fredzannarbor/ar7-climate-assessment}

\paragraph{Disclaimer.}
Not affiliated with or endorsed by the IPCC. All model output is unvalidated.

\bibliographystyle{plainnat}
\bibliography{refs}

\end{document}
